\section{Results}

\subsection{Selected model configuration}

The automated optimization procedure evaluated 100 candidate configurations, of
which 70 completed all four cross-validation folds and 30 were terminated early
by the median pruning rule. Model selection was based on the macro-averaged
area under the precision--recall curve (PR-AUC) at a one-month lead time,
averaged over the four expanding-window folds.

The selected configuration is a Random Forest classifier using 32 of the 52
candidate predictors, achieving a mean cross-validation macro PR-AUC of 0.324
(Table~\ref{tab:best_config}). Tree-based models dominated the upper range of
the search: across completed trials, Random Forest achieved a mean score of
0.305 and XGBoost 0.300, compared with 0.273 for CatBoost, 0.267 for Elastic
Net and 0.263 for Logistic Regression. Fold-level scores for the selected
configuration were 0.433 (validation year 2020), 0.193 (2021), 0.336 (2022) and
0.334 (2023). The comparatively poor performance in 2021 corresponds to a year
without a major drought event, indicating that predictive skill is concentrated
in drought years.

A notable outcome of the predictor search is that the standardized drought
indices (SPI and SPEI) were not retained. The optimization selected the
\textit{auxiliary} meteorological subset, consisting of soil moisture anomalies
and consecutive dry days, in preference to the SPI and SPEI accumulations.
All remaining predictor groups were retained, with the class-lag group
restricted to agricultural impact lags and the elevation group restricted to
mean elevation.

\begin{table}[h]
\centering
\small
\caption{Selected model configuration obtained from the automated optimization procedure.}
\label{tab:best_config}
\begin{tabular}{p{4.5cm} p{9.5cm}}
\toprule
\textbf{Component} & \textbf{Selected value} \\
\midrule
Algorithm & Random Forest \\
Hyperparameters & 400 trees, maximum depth 6, minimum samples per leaf 9, balanced subsample class weighting \\
Number of predictors & 32 (of 52 candidates) \\
Hydro-meteorological & Soil moisture anomaly, consecutive dry days \\
Regional characteristics & Land cover (3), population fraction, mean elevation, soil composition (18) \\
Temporal & Seasonal harmonics, previous impact occurrence, agricultural impact lags (4) \\
Cross-validation macro PR-AUC ($h=1$) & 0.324 \\
\bottomrule
\end{tabular}
\end{table}

\subsection{Forecast performance on the independent test period}

The selected configuration was retrained on the complete development period
(2018--2023) and evaluated on the independent test period (2024--2025), which
comprises 837 region--month observations across 40 NUTS-3 regions.
Table~\ref{tab:test_horizon} reports macro-averaged performance for each
forecast horizon.

Because the impact classes are rare, PR-AUC values must be interpreted relative
to the class prevalence, which represents the expected performance of a
no-skill classifier. The skill score is therefore defined as the difference
between the observed PR-AUC and the prevalence.

At a one-month lead time the model achieves a macro PR-AUC of 0.324 against a
macro prevalence of 0.091, corresponding to a skill of 0.233 and a 3.6-fold
improvement over the no-skill baseline. Performance decreases monotonically
with lead time, from a skill of 0.247 at $h=0$ to 0.182 at $h=3$, while
remaining approximately 2.9 times better than the baseline at the longest
horizon. This gradual decline is consistent with the expectation that
predictive skill decreases with increasing lead time, but indicates that useful
skill is retained across the full zero-to-three-month range rather than
collapsing to random performance.

\begin{table}[h]
\centering
\small
\caption{Macro-averaged test performance by forecast horizon (2024--2025, ten impact classes).}
\label{tab:test_horizon}
\begin{tabular}{lrrrrrrr}
\toprule
\textbf{Horizon} & \textbf{PR-AUC} & \textbf{Prevalence} & \textbf{Skill} & \textbf{ROC-AUC} & \textbf{Precision} & \textbf{Recall} & \textbf{Brier} \\
\midrule
$h=0$ & 0.340 & 0.093 & 0.247 & 0.779 & 0.240 & 0.606 & 0.150 \\
$h=1$ & 0.324 & 0.091 & 0.233 & 0.768 & 0.220 & 0.605 & 0.161 \\
$h=2$ & 0.285 & 0.094 & 0.191 & 0.744 & 0.190 & 0.608 & 0.177 \\
$h=3$ & 0.277 & 0.095 & 0.182 & 0.734 & 0.192 & 0.602 & 0.178 \\
\bottomrule
\end{tabular}
\end{table}

The model operates in a high-recall, low-precision regime, correctly
identifying approximately 60\% of impact occurrences at a precision of
0.19--0.24. This behaviour follows from the balanced class weighting applied
during training and is appropriate for an early-warning context, in which
failing to anticipate an impact is generally more costly than issuing a false
alarm.

\subsection{Performance across impact classes}

Performance varies considerably between impact classes
(Table~\ref{tab:test_class}). Agricultural Economic Loss, Crop Failure and
Yield Reduction, Irrigation Shortage and Water Use Restrictions form a
consistently well-predicted group, with skill scores between 0.315 and 0.363
and ROC-AUC values approaching 0.83. These classes are directly coupled to the
agricultural water balance, which is well represented by the selected soil
moisture and dry-spell predictors.

A second group, comprising Reservoir and Surface Water Shortage, Forest Dieback
and Vegetation Stress, and Freshwater Ecosystem Degradation, is substantially
harder to predict, with skill scores of 0.103--0.143. These impacts depend on
accumulated hydrological deficits and ecological response times that are only
partially captured by the monthly predictors used here.

\begin{table}[h]
\centering
\small
\caption{Test performance by impact class at a one-month lead time ($h=1$), ordered by skill.}
\label{tab:test_class}
\begin{tabular}{p{6.2cm} rrrrr}
\toprule
\textbf{Impact class} & \textbf{Prev.} & \textbf{PR-AUC} & \textbf{Skill} & \textbf{ROC-AUC} & \textbf{Recall} \\
\midrule
Agricultural Economic Loss          & 0.035 & 0.398 & 0.363 & 0.832 & 0.621 \\
Crop Failure \& Yield Reduction     & 0.148 & 0.478 & 0.330 & 0.797 & 0.629 \\
Irrigation Shortage                 & 0.088 & 0.412 & 0.323 & 0.822 & 0.635 \\
Water Use Restrictions              & 0.059 & 0.374 & 0.315 & 0.794 & 0.694 \\
Groundwater Depletion               & 0.109 & 0.377 & 0.268 & 0.752 & 0.637 \\
Wildfire Occurrence                 & 0.029 & 0.220 & 0.191 & 0.799 & 0.625 \\
Water Supply \& Sanitation Issues   & 0.111 & 0.269 & 0.157 & 0.750 & 0.570 \\
Freshwater Ecosystem Degradation    & 0.117 & 0.260 & 0.143 & 0.678 & 0.551 \\
Forest Dieback \& Vegetation Stress & 0.111 & 0.250 & 0.139 & 0.734 & 0.613 \\
Reservoir \& Surface Water Shortage & 0.103 & 0.206 & 0.103 & 0.718 & 0.477 \\
\bottomrule
\end{tabular}
\end{table}

\subsection{Temporal variation in forecast skill}

Forecast skill is strongly seasonal (Figure~\ref{fig:diag_calendar_month}).
Skill is highest during the summer months, reaching 0.490 in August, 0.414 in
July and 0.359 in June, and remains substantial in March (0.371) and April
(0.300). In January, October and November the number of observed impacts is too
small for PR-AUC to be defined, and these months are therefore reported as
missing rather than as zero skill. This reflects the seasonal concentration of
drought impacts in the database rather than a failure of the model.

Skill also depends on the prevailing drought regime
(Table~\ref{tab:drought_bin}). Under dry conditions, defined by an SPEI-12
value below $-1$, the model achieves a macro PR-AUC of 0.440 and a skill of
0.296. Under near-normal conditions skill decreases to 0.224, and under wet
conditions to 0.133. The model therefore performs best precisely in the
conditions under which impact-based drought forecasting is operationally
relevant.

\begin{table}[h]
\centering
\small
\caption{Test performance by SPEI-12 drought regime at $h=1$.}
\label{tab:drought_bin}
\begin{tabular}{lrrrrr}
\toprule
\textbf{Drought regime} & \textbf{Observations} & \textbf{Positives} & \textbf{Prevalence} & \textbf{PR-AUC} & \textbf{Skill} \\
\midrule
Dry         & 1410 & 187 & 0.133 & 0.440 & 0.296 \\
Near-normal & 1650 & 290 & 0.176 & 0.400 & 0.224 \\
Wet         & 5310 & 284 & 0.053 & 0.187 & 0.133 \\
\bottomrule
\end{tabular}
\end{table}

\subsection{Spatial variation in forecast skill}

Regional performance was evaluated for each NUTS-3 region separately
(Figure~\ref{fig:diag_choropleth}). Of the 40 regions, 15 contain sufficient
positive observations in the test period for regional skill to be computed;
for the remaining regions the metric is undefined and these are shown as
missing. This sparsity is a direct consequence of the chosen spatial
resolution: disaggregating a limited number of reported impacts across 40
regions and 21 test months leaves many region--class combinations empty.

Among the regions where skill can be estimated, Groot-Rijnmond attains the
highest value (0.507), followed by Zuidoost-Noord-Brabant (0.397),
Zeeuwsch-Vlaanderen (0.369) and Zuid-Limburg (0.340). The lowest values occur
in Veluwe (0.114), Delft en Westland (0.115) and Utrecht (0.146). The
better-performing regions are predominantly located in the southern and
south-western Netherlands, which correspond to the areas most frequently
affected during the 2018, 2019 and 2022 drought events and therefore provide
the richest training signal.

\subsection{Predictor contributions}

Post-hoc interpretation was performed using SHAP values computed on the test
period at a one-month lead time, averaged across the ten class-specific models
(Figure~\ref{fig:shap}). The dominant predictor is the seasonal cosine
component (mean $|$SHAP$|$ = 0.070), confirming that the annual cycle of
drought impacts is the single strongest signal available to the model. This is
followed by consecutive dry days (0.047), the peat soil fraction (0.034) and
the soil moisture anomaly (0.031).

Two observations follow from this ranking. First, the two retained
hydro-meteorological predictors rank second and fourth overall, confirming
that the optimization procedure discarded the SPI and SPEI accumulations
because they were redundant rather than uninformative. Second, static soil
composition variables occupy several of the leading positions, with peat and
peaty soils contributing most strongly. This supports the inclusion of soil
composition proposed during the exploratory analysis and expert consultation,
and is physically consistent with the elevated drought sensitivity of
organic-rich soils in the Netherlands.

\subsection{Predictability limits across impact classes}

To assess whether the restriction to the ten most frequent impact classes is
justified, a supplementary analysis scored all fifteen candidate classes under
a single shared model configuration using the same cross-validation procedure
(Table~\ref{tab:predictability}). A class was considered predictable if it
achieved a cross-validation skill of at least 0.15 and a PR-AUC of at least
0.25.

Seven classes satisfy both criteria, with skill scores between 0.170 and 0.263.
The remaining eight fall below the threshold, and the lowest-ranked classes
(Inland Waterway Disruption, Broader Economic Disruption, Social Impacts and
Wetland Loss) achieve skill scores below 0.05, indicating essentially no
predictive signal. These classes are both rare and only indirectly related to
hydro-meteorological conditions, and their occurrence in the database is
strongly influenced by media reporting behaviour rather than physical drought
severity.

This result confirms that predictive capability is limited to a subset of
impact types and supports the decision to restrict the prediction target to the
most frequently occurring classes. A second optimization stage restricted to
the predictable classes was also evaluated but did not improve performance
relative to the shared configuration (macro PR-AUC 0.374 versus 0.385), and the
single shared configuration was therefore retained.

\begin{table}[h]
\centering
\small
\caption{Cross-validation predictability of all fifteen candidate impact classes at $h=1$ under a shared model configuration.}
\label{tab:predictability}
\begin{tabular}{p{6.2cm} rrrl}
\toprule
\textbf{Impact class} & \textbf{PR-AUC} & \textbf{Prev.} & \textbf{Skill} & \textbf{Predictable} \\
\midrule
Forest Dieback \& Vegetation Stress & 0.396 & 0.133 & 0.263 & Yes \\
Groundwater Depletion               & 0.459 & 0.212 & 0.247 & Yes \\
Irrigation Shortage                 & 0.333 & 0.101 & 0.232 & Yes \\
Reservoir \& Surface Water Shortage & 0.383 & 0.167 & 0.216 & Yes \\
Water Use Restrictions              & 0.322 & 0.120 & 0.202 & Yes \\
Freshwater Ecosystem Degradation    & 0.415 & 0.215 & 0.200 & Yes \\
Crop Failure \& Yield Reduction     & 0.386 & 0.216 & 0.170 & Yes \\
\midrule
Wildfire Risk Increase              & 0.212 & 0.038 & 0.174 & No \\
Wildfire Occurrence                 & 0.185 & 0.043 & 0.143 & No \\
Agricultural Economic Loss          & 0.218 & 0.101 & 0.117 & No \\
Water Supply \& Sanitation Issues   & 0.243 & 0.157 & 0.086 & No \\
Wetland Loss                        & 0.095 & 0.052 & 0.043 & No \\
Social Impacts                      & 0.094 & 0.053 & 0.042 & No \\
Broader Economic Disruption         & 0.121 & 0.087 & 0.035 & No \\
Inland Waterway Disruption          & 0.064 & 0.038 & 0.026 & No \\
\bottomrule
\end{tabular}
\end{table}